\section{Notion de courant électrique}
Le courant électrique correspond à un déplacement de porteurs de charge. La répartition des courants peut être modélisée de plusieurs manières.
\subsection{Distribution volumique}
$\vv{j}(M,t)$ désigne le vecteur densité volumique de courant électrique au point $M$ et à l'instant $t$.

\formule{Courant en description volumique}
[]
{$$I=\iint \vv{j}.\vv{dS}$$}
[
	\begin{itemize}
		\item $I$ le courant électrique en \unit{A}
		\item $\vv{j}$ le vecteur densité volumique de courant électrique en \unit{A.m^{-2}}
	\end{itemize}
][o]

\subsection{Distribution surfacique}
$\vv{j_S}(M,t)$ désigne le vecteur densité surfacique de courant électrique au point $M$ et à l'instant $t$.

\formule{Courant en description surfacique}
[]
{$$I=\int \vv{j_S}.\vv{dl}$$}
[
	\begin{itemize}
		\item $I$ le courant électrique en \unit{A}
		\item $\vv{j_S}$ le vecteur densité surfacique de courant électrique en \unit{A.m^{-1}}
	\end{itemize}
][o]

\subsection{Distribution linéique}
$I$ désigne le courant dans un objet filiforme.

\section{Propriétés du champ magnétostatique}
\subsection{Équations de Maxwell}

Le comportement du champ magnétique est régi par les équations de Maxwell-Thomson et de Maxwell-Ampère.

\formule{Équation de Maxwell-Thomson}
[]
{$$\divv\vv{B}=0$$}
[
    \begin{itemize}
        \item $\vv{B}$ le champ magnétique (en \unit{T})
    \end{itemize}
][o]

\flashcard{Équation de Maxwell-Thomson}{$$\divv\vv{B}=0$$}

\formule{Équation de Maxwell-Ampère}
[]
{$$\rot\vv{B}=\mu_0\vv{j}+\mu_0\epsilon_0\frac{\partial\vv{E}}{\partial t}$$}
[
    \begin{itemize}
        \item $\vv{B}$ le champ magnétique (en \unit{T})
        \item $\vv{j}$ le vecteur densité volumique de courant électrique (en \unit{A.m^{-2}})
        \item $\vv{E}$ le champ électrique (en \unit{V.m^{-1}})
        \item $\mu_0=4\pi \times \SI{e-7}{H.m^{-1}}$ la perméabilité magnétique du vide
        \item $\epsilon_0=\SI{8.85e-12}{F.m^{-1}}$ la permittivité diélectrique du vide
    \end{itemize}
][o]

\flashcard{Équation de Maxwell-Ampère}{$$\rot\vv{B}=\mu_0\vv{j}+\mu_0\epsilon_0\frac{\partial\vv{E}}{\partial t}$$}

\subsection{Théorème de superposition}
Les équations de Maxwell sont linéaires.

Le champ magnétique résultant de plusieurs distributions de courant est la somme vectorielle des champs magnétiques créés par chacune des distributions de courant.

\subsection{Conservation du flux du champ magnétique}

\formule{Conservation du flux de $\vv{B}$}
[]
{Le champ magnétique est à flux conservatif}
[][o][o]

\flashcard{A quelle condition le champ magnétique est-il à flux conservatif ?}{Aucune, c'est toujours le cas.}

Du fait de cette propriété, l'équation de Maxwell-Thomson est parfois appelée équation de Maxwell-flux.

Un évasement d'un tube de champ s'accompagne donc de la diminution de la norme du champ magnétique.

\colle{Énoncer l'équation de Maxwell, démontrer que le champ magnétique est à flux conservatif et faire le lien avec la topographie des cartes de champ magnétique.}

\subsection{Forces causées par un champ magnétique}

Le champ magnétique exerce une force sur les particules chargées en mouvement. L'expression de la force exercée par le champ magnétique dépend de la description du déplacement des charges.

\subsubsection{Description ponctuelle}

Dans la description ponctuelle, une particule chargée est animée d'une vitesse.

\formule{Force magnétique subie par une particule ponctuelle}
[]
{$$\vv{F}=q\vv{v}\wedge\vv{B}$$}
[
    \begin{itemize}
        \item $\vv{F}$ la partie magnétique de la force de Lorentz subie par la particule (en \unit{N})
        \item $q$ la charge de la particule (en \unit{C})
        \item $\vv{B}$ le champ magnétique (en \unit{T})
        \item $\vv{v}$ la vitesse de la particule chargée (en \unit{C})
    \end{itemize}
][o]

\flashcard{Force de Lorentz}{$$q\vv{E}+q\vv{v}\wedge\vv{B}$$}

\application{Une particule chargée de charge $q$ et de masse $m$ plongée dans un champ manétique uniforme et stationnaire $\vv{B}=B\vv{e_z}$ a une trajectoire circulaire orthogonale à $\vv{B}$. Exprimer sa vitesse angulaire.}

\subsubsection{Desctription linéique}

Dans la description linéique, un fil infiniment fin est parcouru par un courant électrique.

\formule{Force de Laplace}
[]
{$$\delta\vv{F}=I\vv{dl}\wedge\vv{B}$$}
[
    \begin{itemize}
        \item $\delta\vv{F}$ la force magnétique subie par une portion de conducteur (en \unit{N})
        \item $I$ le courant parcourant le conducteur (en \unit{A})
        \item $\vv{dl}$ la longueur de la portion de conducteur orientée dans le même sens que le courant (en \unit{m})
        \item $\vv{B}$ le champ magnétique (en \unit{T})
    \end{itemize}
][o][o]

\flashcard{Force de Laplace}{$$\delta\vv{F}=I\vv{dl}\wedge\vv{B}$$}

\application{Dans les rail de Laplace, une barre traversée par un courant de \SI{5}{A} dirigé selon $\vv{e_y}$ roule sur des rails horizontaux distants de \SI{10}{cm} en présence d'un champ magnétique vertical de \SI{3e-2}{T} dirigé selon $\vv{e_z}$. Calculer la norme de la force magnétique subie par la barre.}

\subsubsection{Description volumique}

Dans la description volumique, le vecteur densité volumique de courant est présent.

\formule{Force de Laplace volumique}
[]
{$$\delta\vv{F}=\vv{j}dV\wedge\vv{B}$$}
[
    \begin{itemize}
        \item $\delta\vv{F}$ la force magnétique (en \unit{N})
        \item $\vv{j}$ le vecteur densité volumique de courant électrique (en \unit{A.m^{-2}})
        \item $dV$ le volume du système (en \unit{m^3})
        \item $\vv{B}$ le champ magnétique (en \unit{T})
    \end{itemize}
][o][o]

\flashcard{Force de Laplace volumique}{$$\delta\vv{F}=\vv{j}dV\wedge\vv{B}$$}

\colle{Citer l'expression de la partie magnétique de la force de Lorentz. En déduire la force de Laplace exercée sur un élément de fil et sur un élément de volume.}

\application{Un fil épais de section \SI{6}{mm^2} et de longueur \SI{10}{m} est parcouru par vecteur densité de courant uniforme $\vv{j}=j\vv{e_x}$ de sorte que le courant totale vaut \SI{10}{A}. Calculer la norme de la force exercée par le champ magnétique terrestre ($\SI{5e-5}{T}\vv{e_z}$), supposé orthogonal au courant.}

\section{Théorème d'Ampère}

\subsection{Invariances du champ magnétique}

Les invariances de la distribution de courant contraignent la forme du champ magnétique.

D'après le principe de Curie, les invariances de la distribution de courant sont aussi des invariances du champ magnétique.

\flashcard{Lien entre les invariances de la distribution de charge et celles du champ magnétiques}{Les invariances de la distribution de courant sont aussi des invariances du champ magnétique (principe de Curie).}

\subsection{Symétries du champ magnétique}

Le champ magnétique est un pseudo vecteur :
\begin{itemize}
    \item le champ magnétique est symétrique par rapport aux plans d'antisymétrie
    \item le champ magnétique est antisymétrique par rapport aux plans de symétrie
\end{itemize}

\formule{Plans de symétrie et champ magnétique}
[]
{Le champ magnétique est orthogonal aux plans de symétrie de la distribution de courant.}
[][o][o]

\flashcard{Lien entre le champ magnétique et les plans de symétrie}{Le champ magnétique est orthogonal aux plans de symétrie de la distribution de courant.}

\application{Déterminer la direction du champ magnétique créé par un fil rectiligne infini et infiniment fin parcourru par un courant.}

\formule{Plans d'antisymétrie et champ magnétique}
[]
{Le champ magnétique est inclus dans les plans d'antisymétrie de la distribution de courant.}
[][o][o]

\flashcard{Lien entre le champ magnétique et les plans d'antisymétrie}{Le champ magnétique est inclus dans les plans d'antisymétrie de la distribution de courant.}

\colle{Énoncer le principe de Curie. Établir le lien existant entre les plans de symétrie et d'antisymétrie de la distribution de courant et la direction du champ magnétique.}

\subsection{Théorème d'Ampère}

Le théorème d'Ampère permet de déterminer le champ magnétique à partir de la distribution de courant.

\formule{Théorème d'Ampère}
[en régime stationnaire]
{$$\oint_{\mathcal{C}}\vv{B}.\vv{dl}=\mu_0 I_\text{enlacé}$$}
[
    \begin{itemize}
        \item $\mathcal{C}$ une courbe fermée orientée
        \item $\vv{B}$ le champ magnétique (en \unit{T})
        \item $\mu_0=4\pi \times \SI{e-7}{H.m^{-1}}$ la perméabilité magnétique du vide
        \item $I_\text{enlacé}$ le courant enlacé par $\mathcal{C}$ (orienté dans le sens direct) (en \unit{A})
    \end{itemize}
][o][o]

\flashcard{Théorème d'Ampère}{$$\oint_{\mathcal{C}}\vv{B}.\vv{dl}=\mu_0 I_\text{enlacé}$$}

\colle{Citer l'équation de Maxwell-Ampère puis établir le théorème d'Ampère.}

\application{Déterminer le champ magnétique créé dans tout l'espace par un fil épais de rayon $R$, rectiligne et infini parcouru par un vecteur densité de courant $\vv{j}$.}
\colle{Déterminer le champ magnétique créé dans tout l'espace par un fil épais de rayon $R$ et infini parcouru par un vecteur densité de courant $\vv{j}$.}

\application{Déterminer le champ magnétique créé dans tout l'espace par un solénoïde infini de rayon $R$ comportant $n$ spires par unité de longueur parcourru par un courant $I$. Le solénoïde est assimilé à une succession de spires circulaires jointives. On suppose le champ magnétique nul à l'extérieur du solénoïde.}
\colle{Déterminer le champ magnétique créé dans tout l'espace par un solénoïde infini de rayon $R$ comportant $n$ spires par unité de longueur parcourru par un courant $I$. Le solénoïde est assimilé à une succession de spires circulaires jointives. On suppose le champ magnétique nul à l'extérieur du solénoïde.}

\application{Déterminer le champ magnétique créé par une bobine torique comportant $N\gg 1$ spires parcourue par un courant $I$.}
\colle{Déterminer le champ magnétique créé par une bobine torique comportant $N\gg 1$ spires parcourue par un courant $I$.}
